\chapter{2010年山东卷（文）}

\section{单选题}

\begin{problem}
已知全集 \(U = \R\)，集合 \(M = \{x \mid x^2 - 4 \leqslant 0\}\)，则 \(\complement_U M =\)（\quad）

\choices
    {\(\{x\mid -2 < x < 2\}\)}
    {\(\{x\mid -2\leqslant x\leqslant 2\}\)}
    {\(\{x \mid x < -2 \text{\ 或\ } x > 2\}\)}
    {\(\{x\mid x\leqslant -2\text{或} x\geqslant 2\}\)}
\end{problem}

\begin{answer}
C
\end{answer}

\begin{solution}
由 \(x^2-4\leqslant0\)，得 \((x-2)(x+2)\leqslant0\)，所以
\(M=\{x\mid -2\leqslant x\leqslant2\}\). 因为全集为 \(\R\)，补集取区间外的部分，故
\(\complement_U M=\{x\mid x<-2\text{ 或 }x>2\}\)，所以选 C.
\end{solution}

\begin{problem}
已知 \(\frac{a + 2\i}{\i} = b + \i\)，其中 \(\i\) 为虚数单位，则 \(a + b =\)（\quad）（\(a\)，\(b \in \R\)）

\choices
    {\(-1\)}
    {\(1\)}
    {\(2\)}
    {\(3\)}
\end{problem}

\begin{answer}
B
\end{answer}

\begin{solution}
因为 \(\dfrac{1}{\i}=-\i\)，所以
\[
\frac{a+2\i}{\i}=a\frac{1}{\i}+2=-a\i+2
\]
与 \(b+\i\) 比较实部和虚部，得 \(b=2\)，\(-a=1\)，即 \(a=-1\). 因此
\(a+b=-1+2=1\)，所以选 B.
\end{solution}

\begin{problem}
函数 \(f(x) = \log_2(3^x + 1)\) 的值域为（\quad）

\choices
    {\((0, +\infty)\)}
    {\([0, +\infty)\)}
    {\((1, +\infty)\)}
    {\([1, +\infty)\)}
\end{problem}

\begin{answer}
A
\end{answer}

\begin{solution}
当 \(x\in\R\) 时，\(3^x\in(0,+\infty)\)，从而 \(3^x+1\in(1,+\infty)\). 由于底数 \(2>1\)，对数函数 \(\log_2 t\) 在 \((0,+\infty)\) 上递增，因此
\(f(x)=\log_2(3^x+1)\in(0,+\infty)\).

反过来，任取 \(u>0\)，令 \(3^x=2^u-1>0\)，则存在实数 \(x\) 使等式成立，且 \(f(x)=u\). 所以值域为 \((0,+\infty)\)，选 A.
\end{solution}

\begin{problem}
在空间，下列命题正确的是（\quad）

\choices
    {平行直线的平行投影重合}
    {平行于同一直线的两个平面平行}
    {垂直于同一平面的两个平面平行}
    {垂直于同一平面的两条直线平行}
\end{problem}

\begin{answer}
D
\end{answer}

\begin{solution}
选项 A 不正确，因为两条平行直线的平行投影可以是两条不同的平行直线，不一定重合.

选项 B 不正确，例如取相交平面 \(x=0\) 与 \(y=0\)，再取直线 \(x=1\)，\(y=1\)（方向为 \(z\) 轴方向）；该直线分别平行于两个平面，但两个平面相交.

选项 C 不正确，例如取底面为 \(xOy\) 平面，则平面 \(x=0\) 与平面 \(y=0\) 都垂直于底面，但它们相交于 \(z\) 轴，不互相平行.

若两条直线都垂直于同一平面，则它们的方向都与该平面的法向量平行，所以两条直线互相平行，选项 D 正确.
\end{solution}

\begin{problem}
设 \(f(x)\) 为定义在 \(\R\) 上的奇函数，当 \(x \geqslant 0\) 时，\(f(x) = 2^x + 2x + b\)（\(b\) 为常数），则 \(f(-1) =\)（\quad）

\choices
    {\(-3\)}
    {\(-1\)}
    {\(1\)}
    {\(3\)}
\end{problem}

\begin{answer}
A
\end{answer}

\begin{solution}
奇函数满足 \(f(0)=0\). 代入题设中 \(x=0\) 的表达式，得 \(1+b=0\)，所以 \(b=-1\).

由题设在非负数上的表达式，\(f(1)=2^1+2\cdot1-1=3\). 又因为 \(f\) 为奇函数，\(f(-1)=-f(1)=-3\)，所以选 A.
\end{solution}

\begin{problem}
在某项体育比赛中，七位裁判为一选手打出的分数如下：\(90\)，\(89\)，\(90\)，\(95\)，\(93\)，\(94\)，\(93\). 去掉一个最高分和一个最低分后，所剩数据的平均分值和方差分别为（\quad）

\choices
    {\(92\)，\(2\)}
    {\(92\)，\(2.8\)}
    {\(93\)，\(2\)}
    {\(93\)，\(2.8\)}
\end{problem}

\begin{answer}
B
\end{answer}

\begin{solution}
去掉最高分 \(95\) 和最低分 \(89\)，剩下的数据为 \(90\)，\(90\)，\(93\)，\(94\)，\(93\). 其平均数为
\[
\bar{x}=\frac{90+90+93+94+93}{5}=92
\]
相对于平均数 \(92\) 的偏差为 \(-2\)，\(-2\)，\(1\)，\(2\)，\(1\)，所以方差为
\[
s^2=\frac{(-2)^2+(-2)^2+1^2+2^2+1^2}{5}=\frac{14}{5}=2.8
\]
因此选 B.
\end{solution}

\begin{problem}
设 \(\{a_n\}\) 是首项大于零的等比数列，则“\(a_1 < a_2\)”是“数列 \(\{a_n\}\) 是递增数列”的（\quad）

\choices
    {充分而不必要条件}
    {必要而不充分条件}
    {充分必要条件}
    {既不充分也不必要条件}
\end{problem}

\begin{answer}
C
\end{answer}

\begin{solution}
设等比数列的公比为 \(q\). 因为 \(a_1>0\)，所以
\[
a_1<a_2=a_1q\quad\Longleftrightarrow\quad q>1
\]
当且仅当 \(q>1\) 时，\(a_{n+1}=a_nq>a_n\) 对所有 \(n\in\N^*\) 成立，即数列为递增数列. 因此“\(a_1<a_2\)”是“数列递增”的充分必要条件，选 C.
\end{solution}

\begin{problem}
已知某生产厂家的年利润 \(y\)（单位：万元）与年产量 \(x\)（单位：万件）的函数关系式为 \(y = -\frac{1}{3}x^{3} + 81x - 234\)，则使该生产厂家获得最大年利润的年产量为（\quad）

\choices
    {\(13\text{ 万件}\)}
    {\(11\text{ 万件}\)}
    {\(9\text{ 万件}\)}
    {\(7\text{ 万件}\)}
\end{problem}

\begin{answer}
C
\end{answer}

\begin{solution}
年产量不能为负数，故取 \(x\geqslant0\). 令
\(y(x)=-\frac{1}{3}x^3+81x-234\)，\(y'(x)=-x^2+81=(9-x)(9+x)\).
当 \(0\leqslant x<9\) 时，\(y'(x)>0\)；当 \(x>9\) 时，\(y'(x)<0\). 因此 \(y(x)\) 在 \([0,9]\) 上递增，在 \([9,+\infty)\) 上递减，最大值在 \(x=9\) 时取得. 所以年产量为 \(9\) 万件，选 C.
\end{solution}

\begin{problem}
已知抛物线 \(y^{2} = 2px\) （\(p > 0\)），过其焦点且斜率为 \(1\) 的直线交抛物线于 \(A\)，\(B\) 两点，若线段 \(AB\) 的中点的纵坐标为 \(2\)，则该抛物线的准线方程为（\quad）

\choices
    {\(x = 1\)}
    {\(x = -1\)}
    {\(x = 2\)}
    {\(x = -2\)}
\end{problem}

\begin{answer}
B
\end{answer}

\begin{solution}
将 \(y^2=2px\) 写成标准形式 \(y^2=4ax\)，得 \(a=\dfrac p2\)，所以焦点为 \(\left(\dfrac p2,0\right)\)，准线为 \(x=-\dfrac p2\).

过焦点且斜率为 \(1\) 的直线为 \(y=x-\dfrac p2\). 代入抛物线方程，得
\[
\left(x-\frac p2\right)^2=2px
\quad\Longleftrightarrow\quad
x^2-3px+\frac{p^2}{4}=0
\]
设交点的横坐标为 \(x_1\)，\(x_2\)，由韦达定理，\(x_1+x_2=3p\). 交点纵坐标分别为 \(x_1-\dfrac p2\)，\(x_2-\dfrac p2\)，故线段 \(AB\) 中点的纵坐标为
\[
\frac{1}{2}\left(x_1-\frac p2+x_2-\frac p2\right)=\frac{x_1+x_2-p}{2}=p
\]
由题意 \(p=2\)，所以准线方程为 \(x=-1\)，选 B.
\end{solution}

\begin{problem}
观察 \((x^{2})^{\prime} = 2x\)，\((x^{4})^{\prime} = 4x^{3}\)，\((\cos x)^{\prime} = -\sin x\)，由归纳推理可得：若定义在 \(\R\) 上的函数 \(f(x)\) 满足 \(f(-x) = f(x)\)，记 \(g(x)\) 为 \(f(x)\) 的导函数，则 \(g(-x) =\)（\quad）

\choices
    {\(f(x)\)}
    {\(-f(x)\)}
    {\(g(x)\)}
    {\(-g(x)\)}
\end{problem}

\begin{answer}
D
\end{answer}

\begin{solution}
对恒等式 \(f(-x)=f(x)\) 两边关于 \(x\) 求导，左边由链式法则得 \(-f'(-x)\)，右边得 \(f'(x)\). 因此
\[
-f'(-x)=f'(x)
\]
记 \(g(x)=f'(x)\)，则 \(g(-x)=-g(x)\)，所以选 D.
\end{solution}

\begin{problem}
函数 \(y = 2^x - x^{2}\) 的图象大致是（\quad）

\choices
    {\choicebitmap[width=0.15\paperwidth]{img/2010/shandong_liberal/q11_optA_nolabel.png}}
    {\choicebitmap[width=0.15\paperwidth]{img/2010/shandong_liberal/q11_optB_nolabel.png}}
    {\choicebitmap[width=0.15\paperwidth]{img/2010/shandong_liberal/q11_optC_nolabel.png}}
    {\choicebitmap[width=0.15\paperwidth]{img/2010/shandong_liberal/q11_optD_nolabel.png}}
\end{problem}

\begin{answer}
A
\end{answer}

\begin{solution}
记 \(F(x)=2^x-x^2\). 当 \(x<0\) 时，\(F'(x)=2^x\ln2-2x>0\)，所以 \(F(x)\) 在 \((-\infty,0)\) 上递增. 又 \(\displaystyle\lim_{x\to-\infty}F(x)=-\infty\)，且 \(F(0)=1>0\)，故在负半轴上恰有一个零点.

当 \(x>0\) 时，\(F''(x)=2^x(\ln2)^2-2\) 严格递增，且 \(F''(0)<0\)，\(F''(3)>0\)，因此 \(F'(x)\) 先减后增. 由 \(F'(0)=\ln2>0\)，\(F'(1)=2\ln2-2<0\)，\(F'(3)=8\ln2-6<0\)，\(F'(4)=16\ln2-8>0\)，可知 \(F'\) 在 \((0,1)\) 和 \((3,4)\) 内各有一个零点，且除此之外没有正零点. 因而 \(F\) 在正半轴上先增、后减、再增.

此外，\(F(1)=1>0\)，\(F(2)=0\)，\(F(3)=-1<0\)，\(F(4)=0\). 结合上述单调性，\(F\) 在 \((0,+\infty)\) 上恰有两个零点 \(x=2\) 和 \(x=4\)，在 \(x=2\) 处由正变负，在 \(x=4\) 处由负变正；再结合负半轴上的一个零点及两端变化，图象与选项 A 一致，所以选 A.
\end{solution}

\begin{problem}
定义平面向量之间的一种运算“\(\odot\)”如下：对任意的 \(\bs{a} = (m, n)\)，\(\bs{b} = (p, q)\)，令 \(\bs{a} \odot \bs{b} = mq - np\). 下面说法错误的是（\quad）

\choices
    {若 \(\bs{a}\) 与 \(\bs{b}\) 共线，则 \(\bs{a} \odot \bs{b} = 0\)}
    {\(\bs{a} \odot \bs{b} = \bs{b} \odot \bs{a}\)}
    {对任意的 \(\lambda \in \R\)，有 \((\lambda \bs{a}) \odot \bs{b} = \lambda (\bs{a} \odot \bs{b})\)}
    {\((\bs{a}\odot \bs{b})^2 +(\bs{a}\cdot \bs{b})^2 = |\bs{a}|^2 |\bs{b}|^2\)}
\end{problem}

\begin{answer}
B
\end{answer}

\begin{solution}
由定义，\(\bs{a}\odot\bs{b}=mq-np\) 是二维行列式. 若 \(\bs{a}\) 与 \(\bs{b}\) 共线，则 \(mq-np=0\)，所以选项 A 正确.

交换两个向量后，
\[
\bs{b}\odot\bs{a}=pn-qm=-(mq-np)=-(\bs{a}\odot\bs{b})
\]
一般不等于 \(\bs{a}\odot\bs{b}\)，故选项 B 错误.

对任意 \(\lambda\in\R\)，
\[
(\lambda\bs{a})\odot\bs{b}=(\lambda m)q-(\lambda n)p=\lambda(mq-np)=\lambda(\bs{a}\odot\bs{b})
\]
所以选项 C 正确. 另外
\[
(mq-np)^2+(mp+nq)^2=(m^2+n^2)(p^2+q^2)=|\bs{a}|^2|\bs{b}|^2
\]
所以选项 D 也正确. 因此错误的是 B.
\end{solution}

\section{填空题}

\begin{problem}
执行如图所示的程序框图，若输入 \(x = 4\)，则输出 \(y\) 的值为\fillinblank{}.

\texfigure[max width=0.36\linewidth]{img_repaint/2010/shandong_liberal/q13_fig1.tex}
\end{problem}

\begin{answer}
\(-\frac{5}{4}\)
\end{answer}

\begin{solution}
输入初值为 \(x_0=4\). 每轮先计算 \(y=\dfrac12x-1\)，再判断 \(|y-x|<1\)，若不满足则令 \(x=y\) 后继续.

第一轮：\(y=\dfrac12\cdot4-1=1\)，且 \(|1-4|=3\geqslant1\)，故令 \(x=1\).

第二轮：\(y=\dfrac12\cdot1-1=-\dfrac12\)，且 \(\left|-\dfrac12-1\right|=\dfrac32\geqslant1\)，故令 \(x=-\dfrac12\).

第三轮：\(y=\dfrac12\cdot\left(-\dfrac12\right)-1=-\dfrac54\)，且 \(\left|-\dfrac54+\dfrac12\right|=\dfrac34<1\)，满足判断条件，输出 \(y=-\dfrac54\).
\end{solution}

\begin{problem}
已知 \(x\)，\(y \in \R_+\)，且满足 \(\frac{x}{3} + \frac{y}{4} = 1\)，则 \(xy\) 的最大值为\fillinblank{}.
\end{problem}

\begin{answer}
\(3\)
\end{answer}

\begin{solution}
令 \(u=\dfrac{x}{3}\)，\(v=\dfrac{y}{4}\)，则 \(u\geqslant0\)，\(v\geqslant0\)，且 \(u+v=1\). 由基本不等式，
\[
uv\leqslant\left(\frac{u+v}{2}\right)^2=\frac14
\]
因此
\[
xy=(3u)(4v)=12uv\leqslant3
\]
当且仅当 \(u=v=\dfrac12\)，即 \(x=\dfrac32\)，\(y=2\) 时等号成立，所以 \(xy\) 的最大值为 \(3\).
\end{solution}

\begin{problem}
在 \(\triangle ABC\) 中，\(A\)，\(B\)，\(C\) 所对的边分别为 \(a\)，\(b\)，\(c\). 若 \(a = \sqrt{2}\)，\(b = 2\)，\(\sin B + \cos B = \sqrt{2}\)，则角 \(A\) 的大小为\fillinblank{}.
\end{problem}

\begin{answer}
\(\frac{\pi}{6}\)
\end{answer}

\begin{solution}
因为 \(B\in(0,\pi)\)，所以
\[
\sin B+\cos B=\sqrt2\sin\left(B+\frac\pi4\right)=\sqrt2
\]
此时 \(B+\dfrac\pi4\in\left(\dfrac\pi4,\dfrac{5\pi}{4}\right)\)，在该区间内正弦值为 \(1\) 只可能取在 \(\dfrac\pi2\)，故 \(B=\dfrac\pi4\).

由正弦定理，
\(\frac{a}{\sin A}=\frac{b}{\sin B}\)，\(\Longrightarrow\)，\(\sin A=\frac{a\sin B}{b}=\frac{\sqrt2\cdot\frac{\sqrt2}{2}}{2}=\frac12\).
所以 \(A=\dfrac\pi6\) 或 \(A=\dfrac{5\pi}{6}\). 但三角形内角满足 \(A+B<\pi\)，而 \(B=\dfrac\pi4\)，故 \(A<\dfrac{3\pi}{4}\)，排除 \(\dfrac{5\pi}{6}\)，得到 \(A=\dfrac\pi6\).
\end{solution}

\begin{problem}
已知圆 \(C\) 过点 \((1,0)\)，且圆心在 \(x\) 轴的正半轴上，直线 \(l: y = x - 1\) 被圆 \(C\) 所截得的弦长为 \(2\sqrt{2}\)，则过圆心且与直线 \(l\) 垂直的直线的方程为\fillinblank{}.
\end{problem}

\begin{answer}
\(x+y-3=0\)
\end{answer}

\begin{solution}
设圆心为 \(O=(r,0)\)，其中 \(r>0\). 因为圆过点 \((1,0)\)，所以半径为 \(R=|r-1|\). 直线 \(l\) 写成 \(x-y-1=0\)，圆心到直线的距离为
\[
d=\frac{|r-1|}{\sqrt2}=\frac{R}{\sqrt2}
\]
弦长公式给出
\(2\sqrt{R^2-d^2}=2\sqrt2\)，\(\Longrightarrow\)，\(2\sqrt{R^2-\frac{R^2}{2}}=2\sqrt2\)，\(\Longrightarrow\)，\(R=2\).
于是 \(|r-1|=2\)，得 \(r=3\) 或 \(r=-1\). 由 \(r>0\)，取 \(r=3\)，所以圆心为 \((3,0)\).

直线 \(l\) 的斜率为 \(1\)，过圆心且垂直于 \(l\) 的直线斜率为 \(-1\)，方程为 \(y=-x+3\)，即 \(x+y-3=0\).
\end{solution}

\section{解答题}

\begin{problem}
已知函数 \(f(x) = \sin(\pi - \omega x)\cos \omega x + \cos^2\omega x\) （\(\omega > 0\)） 的最小正周期为 \(\pi\).

\begin{enumerate}
\item 求 \(\omega\) 的值；
\item 将函数 \(y = f(x)\) 的图象上各点的横坐标缩短到原来的 \(\frac{1}{2}\)，纵坐标不变，得到函数 \(y = g(x)\) 的图象，求函数 \(y = g(x)\) 在区间 \(\left[0, \frac{\pi}{16}\right]\) 上的最小值.
\end{enumerate}
\end{problem}

\begin{solution}
\begin{enumerate}
\item 由 \(\sin(\pi-\omega x)=\sin\omega x\)，得

\[
f(x)=\sin\omega x\cos\omega x+\cos^2\omega x=\frac{1}{2}+\frac{1}{2}\left(\sin2\omega x+\cos2\omega x\right)=\frac{1}{2}+\frac{\sqrt{2}}{2}\sin\left(2\omega x+\frac{\pi}{4}\right)
\]

由于 \(\omega>0\)，函数 \(f(x)\) 的最小正周期为 \(\frac{\pi}{\omega}\). 由题意 \(\frac{\pi}{\omega}=\pi\)，所以 \(\omega=1\).

\item 将函数 \(y=f(x)\) 的图象上各点的横坐标缩短到原来的 \(\frac{1}{2}\)，相当于用 \(2x\) 代替原函数中的 \(x\)，故

\[
g(x)=f(2x)=\frac{1}{2}+\frac{\sqrt{2}}{2}\sin\left(4x+\frac{\pi}{4}\right)
\]

当 \(x\in\left[0,\frac{\pi}{16}\right]\) 时，\(4x+\frac{\pi}{4}\in\left[\frac{\pi}{4},\frac{\pi}{2}\right]\)，正弦函数在该区间上单调递增，因此

\[
g(x)\geqslant \frac{1}{2}+\frac{\sqrt{2}}{2}\sin\frac{\pi}{4}=1
\]

当 \(x=0\) 时等号成立，所以函数 \(y=g(x)\) 在所给区间上的最小值为 \(1\).
\end{enumerate}
\end{solution}

\begin{problem}
已知等差数列 \(\{a_n\}\) 满足：\(a_3 = 7\)，\(a_5 + a_7 = 26\). 数列 \(\{a_n\}\) 的前 \(n\) 项和为 \(S_n\).

\begin{enumerate}
\item 求 \(a_n\) 及 \(S_n\)；
\item 令 \(b_n = \frac{1}{a_n^2 - 1}\) （\(n \in \N^*\)），求数列 \(\{b_n\}\) 的前 \(n\) 项和 \(T_n\).
\end{enumerate}
\end{problem}

\begin{solution}
\begin{enumerate}
\item 设等差数列 \(\{a_n\}\) 的首项为 \(a_1\)，公差为 \(d\)，则

\[
\begin{cases}
a_1+2d=7,\\
(a_1+4d)+(a_1+6d)=26.
\end{cases}
\]

即 \(a_1+5d=13\). 两式相减得 \(3d=6\)，所以 \(d=2\)，进而 \(a_1=3\). 因此

\[
a_n=a_1+(n-1)d=3+2(n-1)=2n+1
\]

数列 \(\{a_n\}\) 的前 \(n\) 项和为

\[
S_n=\frac{n(a_1+a_n)}{2}=\frac{n(3+2n+1)}{2}=n(n+2)
\]

\item 因为 \(a_n=2n+1\)，所以

\[
b_n=\frac{1}{a_n^2-1}=\frac{1}{(2n+1)^2-1}=\frac{1}{4n(n+1)}=\frac{1}{4}\left(\frac{1}{n}-\frac{1}{n+1}\right)
\]

于是

\[
T_n=\sum_{k=1}^{n}b_k=\frac{1}{4}\sum_{k=1}^{n}\left(\frac{1}{k}-\frac{1}{k+1}\right)=\frac{1}{4}\left(1-\frac{1}{n+1}\right)=\frac{n}{4(n+1)}
\]
\end{enumerate}
\end{solution}

\begin{problem}
一个袋中装有四个形状大小完全相同的球，球的编号分别为 \(1\)，\(2\)，\(3\)，\(4\).

\begin{enumerate}
\item 从袋中随机取两个球，求取出的球的编号之和不大于 \(4\) 的概率；
\item 先从袋中随机取一个球，该球的编号为 \(m\)，将球放回袋中，然后再从袋中随机取一个球，该球的编号为 \(n\)，求 \(n < m + 2\) 的概率.
\end{enumerate}
\end{problem}

\begin{solution}
\begin{enumerate}
\item 从四个球中随机取两个球，所有可能的取法共有 \(\mathrm C_4^2=6\) 种，且等可能. 编号之和不大于 \(4\) 的取法为 \(\{1,2\}\) 和 \(\{1,3\}\)，共有 \(2\) 种，所以所求概率为

\[
P=\frac{2}{6}=\frac{1}{3}
\]

\item 取球后放回，所以 \(m\)，\(n\) 的所有可能结果为 \(16\) 个有序数对，且等可能. 当 \(m=1\) 时，满足 \(n<m+2\) 的 \(n\) 有 \(1\)，\(2\)，共 \(2\) 个；当 \(m=2\) 时有 \(3\) 个；当 \(m=3\) 或 \(m=4\) 时均有 \(4\) 个.

因此满足条件的有序数对共有 \(2+3+4+4=13\) 个，所求概率为

\[
P=\frac{13}{16}
\]
\end{enumerate}
\end{solution}

\begin{problem}
在如图所示的几何体中，四边形 \(ABCD\) 是正方形，\(MA \perp\) 平面 \(ABCD\)，\(PD \parallel MA\)，\(E\)，\(G\)，\(F\) 分别为 \(MB\)，\(PB\)，\(PC\) 的中点，且 \(AD = PD = 2MA\).

\begin{enumerate}
\item 求证：平面 \(EFG\) \(\perp\) 平面 \(PDC\)；
\item 求三棱锥 \(P - MAB\) 与四棱锥 \(P - ABCD\) 的体积之比.
\end{enumerate}

\bitmapfigure[max width=0.55\linewidth]{img/2010/shandong_liberal/q20_fig1.png}
\end{problem}

\begin{solution}
设 \(MA=h\)，建立空间直角坐标系，使 \(A\) 为原点，\(AB\)，\(AD\)，\(AM\) 分别为 \(x\)，\(y\)，\(z\) 轴的正方向. 由题意可取

\(A=(0,0,0)\)，\(B=(2h,0,0)\)，\(C=(2h,2h,0)\)，\(D=(0,2h,0)\).

\(M=(0,0,h)\)，\(P=(0,2h,2h)\).

因为 \(E\)，\(G\)，\(F\) 分别为 \(MB\)，\(PB\)，\(PC\) 的中点，所以

\(E=\left(h,0,\frac{h}{2}\right)\)，\(G=(h,h,h)\)，\(F=(h,2h,h)\).

三点 \(E\)，\(G\)，\(F\) 的横坐标均为 \(h\)，且 \(\overrightarrow{EG}=(0,h,h/2)\)，\(\overrightarrow{EF}=(0,2h,h/2)\) 不成比例，故平面 \(EFG\) 的方程为 \(x=h\). 三点 \(P\)，\(D\)，\(C\) 的纵坐标均为 \(2h\)，且 \(\overrightarrow{PD}=(0,0,-2h)\)，\(\overrightarrow{PC}=(2h,0,-2h)\) 不成比例，故平面 \(PDC\) 的方程为 \(y=2h\). 这两个平面的一个法向量分别为 \((1,0,0)\) 和 \((0,1,0)\)，它们的内积为 \(0\)，所以

\[
\text{平面 }EFG\perp\text{平面 }PDC
\]

三角形 \(MAB\) 的面积为 \(\frac{1}{2}\cdot 2h\cdot h=h^2\)，点 \(P\) 到平面 \(MAB\)（即 \(y=0\)）的距离为 \(2h\)，所以

\[
V_{P-MAB}=\frac{1}{3}h^2\cdot 2h=\frac{2}{3}h^3
\]

四边形 \(ABCD\) 的面积为 \(4h^2\)，点 \(P\) 到平面 \(ABCD\)（即 \(z=0\)）的距离为 \(2h\)，所以

\[
V_{P-ABCD}=\frac{1}{3}\cdot 4h^2\cdot 2h=\frac{8}{3}h^3
\]

因此

\[
V_{P-MAB}:V_{P-ABCD}=\frac{2}{3}h^3:\frac{8}{3}h^3=1:4
\]
\end{solution}

\begin{problem}
已知函数 \(f(x) = \ln x - ax + \frac{1 - a}{x} - 1\) （\(a \in \R\)）.

\begin{enumerate}
\item 当 \(a = -1\) 时，求曲线 \(y = f(x)\) 在点 \((2, f(2))\) 处的切线方程；
\item 当 \(a \leqslant \frac{1}{2}\) 时，讨论 \(f(x)\) 的单调性；
\end{enumerate}
\end{problem}

\begin{solution}
\begin{enumerate}
\item 当 \(a=-1\) 时，函数为 \(f(x)=\ln x+x+\frac{2}{x}-1\)，其定义域为 \(x>0\). 因为

\[
f'(x)=\frac{1}{x}+1-\frac{2}{x^2}=\frac{(x-1)(x+2)}{x^2}
\]

所以 \(f(2)=\ln2+2\)，\(f'(2)=1\). 曲线在点 \((2,f(2))\) 处的切线方程为

\[
y-(\ln2+2)=x-2
\]

即 \(y=x+\ln2\).

\item 函数的定义域为 \(x>0\)，且

\[
f'(x)=\frac{1}{x}-a-\frac{1-a}{x^2}=\frac{(x-1)(1-a-ax)}{x^2}
\]

当 \(a\leqslant0\) 时，\(1-a-ax>0\)，所以 \(f'(x)\) 在 \((0,1)\) 上小于 \(0\)，在 \((1,+\infty)\) 上大于 \(0\). 因此 \(f(x)\) 在 \((0,1)\) 上单调递减，在 \((1,+\infty)\) 上单调递增.

当 \(0<a<\frac{1}{2}\) 时，令 \(1-a-ax=0\)，得 \(x=\frac{1-a}{a}>1\). 因此 \(f'(x)\) 的符号依次为负、正、负，故 \(f(x)\) 在 \(\left(0,1\right)\) 上单调递减，在 \(\left(1,\frac{1-a}{a}\right)\) 上单调递增，在 \(\left(\frac{1-a}{a},+\infty\right)\) 上单调递减.

当 \(a=\frac{1}{2}\) 时，

\[
f'(x)=-\frac{(x-1)^2}{2x^2}\leqslant0
\]

所以 \(f(x)\) 在 \((0,+\infty)\) 上单调递减.
\end{enumerate}
\end{solution}

\begin{problem}
如图，已知椭圆 \(\frac{x^2}{a^2} + \frac{y^2}{b^2} = 1\) （\(a > b > 0\)） 过点 \(\left(1, \frac{\sqrt{2}}{2}\right)\)，离心率为 \(\frac{\sqrt{2}}{2}\)，左、右焦点分别为 \(F_1\)，\(F_2\). 点 \(P\) 为直线 \(l: x + y = 2\) 上且不在 \(x\) 轴上的任意一点，直线 \(PF_1\) 和 \(PF_2\) 与椭圆的交点分别为 \(A\)，\(B\) 和 \(C\)，\(D\)，\(O\) 为坐标原点.

\begin{enumerate}
\item 求椭圆的标准方程；
\item 设直线 \(PF_1\)，\(PF_2\) 的斜率分别为 \(k_1\)，\(k_2\).

\begin{enumerate}
\item 证明：\(\frac{1}{k_1} - \frac{3}{k_2} = 2\)；
\item 问直线 \(l\) 上是否存在点 \(P\)，使得直线 \(OA\)，\(OB\)，\(OC\)，\(OD\) 的斜率 \(k_{OA}\)，\(k_{OB}\)，\(k_{OC}\)，\(k_{OD}\) 满足 \(k_{OA} + k_{OB} + k_{OC} + k_{OD} = 0\)？ 若存在，求出所有满足条件的点 \(P\) 的坐标；若不存在，说明理由.
\end{enumerate}
\end{enumerate}

\bitmapfigure[max width=0.45\linewidth]{img/2010/shandong_liberal/q22_fig1.png}
\end{problem}

\begin{solution}
\begin{enumerate}
\item 由椭圆的离心率得

\[
\frac{c^2}{a^2}=\frac{1}{2}
\]

又因为 \(c^2=a^2-b^2\)，所以 \(b^2=\frac{a^2}{2}\). 又因为椭圆过点 \(\left(1,\frac{\sqrt{2}}{2}\right)\)，所以

\[
\frac{1}{a^2}+\frac{\frac{1}{2}}{b^2}=1
\]

代入 \(b^2=\frac{a^2}{2}\)，得 \(\frac{2}{a^2}=1\)，即 \(a^2=2\)，从而 \(b^2=1\). 因此椭圆的标准方程为

\[
\frac{x^2}{2}+y^2=1
\]

\item 由上式可知焦点为 \(F_1=(-1,0)\)，\(F_2=(1,0)\). 设 \(P=(t,2-t)\)，则 \(t\ne2\). 当 \(t\ne-1\)，\(1\)，\(2\) 时，直线 \(PF_1\)，\(PF_2\) 的斜率有定义，且

\(k_1=\frac{2-t}{t+1}\)，\(k_2=\frac{2-t}{t-1}\).

\begin{enumerate}
\item

\[
\frac{1}{k_1}-\frac{3}{k_2}=\frac{t+1}{2-t}-\frac{3(t-1)}{2-t}=2
\]

\item 先设 \(t\ne-1\)，\(1\)，\(2\). 设直线 \(PF_1\) 与椭圆的交点 \(A\)，\(B\) 的横坐标分别为 \(x_A\)，\(x_B\). 直线 \(PF_1\) 的方程为 \(y=k_1(x+1)\)，代入椭圆方程得

\[
(1+2k_1^2)x^2+4k_1^2x+2k_1^2-2=0
\]

由韦达定理，\(x_A+x_B=-\frac{4k_1^2}{1+2k_1^2}\)，\(x_Ax_B=\frac{2(k_1^2-1)}{1+2k_1^2}\).

又因为 \(k_{OA}=\frac{y_A}{x_A}=k_1\left(1+\frac{1}{x_A}\right)\)，所以

\[
k_{OA}+k_{OB}=k_1\left(2+\frac{x_A+x_B}{x_Ax_B}\right)=\frac{2k_1}{1-k_1^2}
\]

同理，直线 \(PF_2\) 的方程为 \(y=k_2(x-1)\). 设其与椭圆交点 \(C\)，\(D\) 的横坐标分别为 \(x_C\)，\(x_D\)，则

\[
(1+2k_2^2)x^2-4k_2^2x+2k_2^2-2=0
\]

从而 \(x_C+x_D=\frac{4k_2^2}{1+2k_2^2}\)，\(x_Cx_D=\frac{2(k_2^2-1)}{1+2k_2^2}\).

因为 \(k_{OC}=k_2\left(1-\frac{1}{x_C}\right)\)，所以

\[
k_{OC}+k_{OD}=k_2\left(2-\frac{x_C+x_D}{x_Cx_D}\right)=\frac{2k_2}{1-k_2^2}
\]

因此所求条件等价于

\[
\frac{2k_1}{1-k_1^2}+\frac{2k_2}{1-k_2^2}=0
\]

将 \(k_1=\frac{2-t}{t+1}\)，\(k_2=\frac{2-t}{t-1}\) 代入，得

\[
2(2-t)\left(\frac{t+1}{3(2t-1)}+\frac{t-1}{2t-3}\right)=0
\]

由于题目要求四条直线 \(OA\)，\(OB\)，\(OC\)，\(OD\) 的斜率都存在，所以 \(x_Ax_Bx_Cx_D\ne0\). 由上面的积式，须有 \(k_1^2\ne1\)，\(k_2^2\ne1\)，即 \(t\ne\frac{1}{2}\)，\(\frac{3}{2}\). 因而在当前情形下各分母非零，所以

\[
(t+1)(2t-3)+3(t-1)(2t-1)=2t(4t-5)=0
\]

故 \(t=0\) 或 \(t=\frac{5}{4}\). 此时相应的点分别为 \(P=(0,2)\) 或 \(P=\left(\frac{5}{4},\frac{3}{4}\right)\).

这两点均不在 \(x\) 轴上，且对应的四条直线 \(OA\)，\(OB\)，\(OC\)，\(OD\) 的斜率都有定义，代入上式可知均满足条件.

还需检查 \(t=-1\) 和 \(t=1\) 两个未包含在上述斜率参数化中的情形. 当 \(t=-1\) 时，\(P=(-1,3)\)，直线 \(PF_1\) 为 \(x=-1\)，其与椭圆的交点为 \(\left(-1,\frac{\sqrt{2}}{2}\right)\)，\(\left(-1,-\frac{\sqrt{2}}{2}\right)\)，故 \(k_{OA}+k_{OB}=0\). 此时 \(k_2=-\frac{3}{2}\)，由上面对 \(PF_2\) 的韦达计算式得 \(k_{OC}+k_{OD}=\frac{2k_2}{1-k_2^2}=\frac{12}{5}\ne0\)，不满足条件. 当 \(t=1\) 时，\(P=(1,1)\)，直线 \(PF_2\) 为 \(x=1\)，故 \(k_{OC}+k_{OD}=0\)；此时 \(k_1=\frac{1}{2}\)，由上面对应的韦达计算式得 \(k_{OA}+k_{OB}=\frac{2k_1}{1-k_1^2}=\frac{4}{3}\ne0\)，也不满足条件. 因此满足条件的点 \(P\) 为 \((0,2)\) 和 \(\left(\frac{5}{4},\frac{3}{4}\right)\).
\end{enumerate}
\end{enumerate}
\end{solution}
